\documentclass[11pt]{article}

% Change "review" to "final" to generate the final (sometimes called camera-ready) version.
\usepackage[review]{acl}

% Standard package includes
\usepackage{times}
\usepackage{latexsym}
\usepackage[T1]{fontenc}
\usepackage[utf8]{inputenc}
\usepackage{microtype}
\usepackage{inconsolata}
\usepackage{graphicx}
\usepackage{booktabs}
\usepackage{algorithm}
\usepackage{algpseudocode}

\title{Beyond Aggregate Accuracy: Recovery and Regression Tracking in Reflective Prompt Optimization}

% ARR review submissions must be fully anonymous. Replace with real names/
% affiliations only in the camera-ready (\usepackage[final]{acl}) version.
\author{Anonymous ACL submission}

\begin{document}
\maketitle

\begin{abstract}
Automatic prompt optimization (APO) methods that use an LLM to iteratively
rewrite a failing prompt are almost universally evaluated by a single
before/after accuracy delta per benchmark. This aggregate number cannot
distinguish an edit that fixes several failures while breaking none from one
that fixes the same failures while silently breaking others, and it does not
verify that gains on the data used to guide editing transfer to unseen
examples. We introduce \textbf{Reflective FDPO}, a feedback-driven prompt
optimization mechanism that, every round, computes exactly which examples
were recovered and which were newly broken relative to the previous round,
on both the optimizer's own working set and a disjoint validation set, and
feeds this per-item outcome back into the optimizer's next edit.
Across five benchmarks (LegalBench-Hearsay, MMLU, IFEval/IFBench, AIME, and
PUPA) and two solver models (GPT-4o-mini and GPT-4.1 Mini), this fine-grained
tracking surfaces a failure mode that aggregate accuracy alone would miss:
validation accuracy can rise substantially while a sealed, disjoint test set
regresses, with no literal memorization involved. The same per-round
instrumentation let us catch and fix two concrete mechanism bugs mid-project,
each with a documented before/after effect. We report results transparently,
including benchmarks where the mechanism helps one solver model and hurts
another, and release the full recovery/regression trace for every reported
run.
\end{abstract}

\section{Introduction}
\label{sec:intro}

The growing use of LLM-powered applications has made prompt design a
first-class engineering problem: small changes in wording can substantially
shift task accuracy, and hand-tuning prompts by trial and error does not
scale across tasks or models. This has motivated automatic prompt
optimization (APO), in which an LLM iteratively rewrites a task prompt using
its own generated feedback on failures, in place of gradient-based updates.
Across this line of work, almost every method is still evaluated the same
way: a single before/after aggregate accuracy delta per benchmark. We argue
this shared evaluation practice obscures something that matters for anyone
trying to use or build on these methods: the aggregate delta cannot
distinguish an edit that fixed several failures while breaking none from one
that fixed the same failures while silently breaking others, and it does not
verify that gains measured on the data used to guide editing transfer to
genuinely unseen examples.

Much as a human prompt engineer reviewing an edit would first check what it
broke and what it fixed before making another change, we argue the optimizer
itself should be shown this same per-item outcome (except on the held-out
test set), round over round, rather than only a final score. We build
\textbf{Reflective FDPO} (Feedback-Driven Prompt Optimization), a mechanism
that, every round, computes exactly which examples were recovered and which
were newly broken relative to the previous round, on both the optimizer's
own working set and a disjoint validation set, and feeds this per-item
outcome directly back into the optimizer's next edit.

We evaluate Reflective FDPO across five benchmarks: LegalBench-Hearsay
(legal classification), MMLU (multi-subject multiple-choice),
IFEval/IFBench (verifiable instruction-following), AIME (competition
mathematics), and PUPA (privacy-conscious delegation), using small,
inexpensive solver models (GPT-4o-mini and GPT-4.1 Mini) rather than the
mid-scale models used by directly comparable prior work, and we are explicit
throughout about where our numbers are, and are not, comparable to those
prior results. Making per-round regression visible this way surfaces a
failure mode that a purely aggregate evaluation would miss entirely:
\textbf{validation and test accuracy can diverge even without memorized
wording}. We document this directly, and distinguish it explicitly from a
separate, previously-observed failure mode in which the optimizer copies
literal wording from the examples it is shown (\S\ref{sec:method}).

\subsection{Research questions}

We organize the rest of the paper around four research questions. This
framing is being actively refined as new evidence comes in; we state plainly
which questions are answered by current evidence and which remain open.

\begin{itemize}
\item \textbf{RQ1 (mechanism value):} does exposing the optimizer to the
measured outcome of its own previous edit change optimization behavior
relative to failure-only feedback? \emph{Partially answered.}
\S\ref{sec:results} documents concrete mechanism revisions made specifically
because outcome information was being discarded.
\item \textbf{RQ2 (generalization gap):} can improvement on the data used to
guide editing diverge from improvement on genuinely unseen data, and can
per-round regression tracking detect this as it happens? \emph{Answered.}
\S\ref{sec:results} documents concrete cases where validation rose
substantially while sealed test fell.
\item \textbf{RQ3 (mechanism attribution):} which specific components of the
mechanism are responsible for observed behavior changes? \emph{Partially
answered.}
\item \textbf{RQ4 (generality across solver families):} does this transfer
to open-weight models at a scale comparable to prior work (Qwen, Llama)?
\emph{Open, in progress.}
\end{itemize}

\subsection{Contributions}

\begin{enumerate}
\item \textbf{Reflective FDPO}, an outcome-aware prompt-optimization
mechanism that, every round, shows the optimizer exactly which items its own
previous edit recovered and regressed, on both its working set and a
disjoint validation set.
\item \textbf{A direct empirical demonstration of a validation/test
generalization gap}, with the optimizer's actual edit text inspected and
memorization ruled out.
\item \textbf{Two mechanism revisions made in direct, documented response to
observed evidence}, each reported with the specific case that motivated it.
\item \textbf{A transparent, model-dependent result}: the same mechanism, on
the same benchmark, is a net regression for one solver model and a
substantial improvement for another once two concrete bugs were fixed.
\item A reusable evaluation formalism (per-round recovered/regressed
tracking) applicable to any prompt-optimization method, not only ours.
\end{enumerate}

\section{Related Work}
\label{sec:related}

Automatic prompt optimization has progressed from token-level, discrete
search~\citep{shin2020autoprompt}, to textual-gradient beam search over
whole-prompt edits~\citep{pryzant2023protegi,yuksekgonul2025textgrad}, to
evolutionary and population-based search~\citep{guo2024evoprompt,
fernando2024promptbreeder}, to structured multi-stage
pipelines~\citep{opsahlong2024miprov2,schnabel2024sammo}, and most recently
to reflective, ``LLM-as-optimizer'' approaches that treat the model's own
critique of its failures as the search signal. The most directly comparable
recent work is GEPA~\citep{agrawal2026gepa}, which combines reflective
prompt mutation with Pareto-frontier candidate selection over an
evolutionary pool of prompts, and reports results on AIME, IFBench, and PUPA
using GPT-4.1 Mini among other models: three benchmarks this paper also
uses, under different train/test protocols and mechanisms (single-lineage,
whole-prompt rewrite here vs.\ Pareto-based multi-candidate evolution
there). Unlike GEPA and prior APO work generally, we do not report only an
aggregate accuracy delta: every result in this paper is accompanied by the
specific set of items recovered and regressed to produce it.

A closely related line of work applies clustered-error refinement with a
whole-document regression gate to a legal classification
task~\citep{zha2026trace2policy}, reporting a gain from 79.7\% to 93.8\%
accuracy on LegalBench-Hearsay for one of six tested solver models. That
paper's own appendix reports one refinement round partly diagnosed using the
nominally held-out test set, which is not a fully sealed evaluation. We
investigated whether this class of result is sensitive to test-set exposure
by running our own mechanism with the sealed test pool deliberately
substituted as the mining set, a diagnostic condition we flag as invalid for
any real accuracy claim. This produced a gain from 68.8\% to 95.3\%
(+26.6pp) on the same dataset, a larger jump than the published result,
using an unrelated and much simpler mechanism. We report this as a
methodological caution, not an accusation: it demonstrates that test-set
exposure alone can produce gains of this magnitude regardless of the
underlying refinement mechanism, and that a genuinely sealed test set, used
for every other result in this paper, is necessary to distinguish real
generalization from this effect.

\section{Method}
\label{sec:method}

\subsection{Problem formulation}

Following the paper-faithful formalization we build on, prompt optimization
is $p_{\text{new}} = \textsc{LLMOptimize}(p_{\text{old}}, E_{\text{fail}},
E_{\text{gold}})$: an optimizer LLM, given the current prompt, a set of
solver failures $E_{\text{fail}}$, and a set of currently-correct examples
$E_{\text{gold}}$ (to avoid breaking what already works), returns a full
rewritten prompt. Every prompt in our schema has five fixed sections: System
Role, Context, Task Details, Constraints, Output Format, so structural edits
are always well-formed and comparable across rounds.
Solver, optimizer, and judge are always distinct model roles, so a solver is
never also its own optimizer or judge.

\subsection{Mechanism evolution}

The final mechanism is the product of three earlier designs, each abandoned
for a concrete, measured reason. An early section-local design, in which an
LLM judge attributed each failure to one prompt section before a scoped edit
was attempted, was retired because judge attribution was noisy and a strict
regression gate rejected over 40\% of proposed edits with no clear
improvement trend. A subsequent whole-prompt design using exact-substring
find/replace edits was retired after three measured failures: a mean effect
statistically indistinguishable from zero across repeated seeds, a
regression gate that rejected net-positive trades outright, and
round-to-round oscillation rather than compounding improvement. One
intermediate design, which explicitly encouraged the optimizer to include
``worked examples'' in the rewritten prompt, caused the optimizer to paste
training cases into the prompt nearly verbatim; test accuracy crashed by 6.8
percentage points before this was diagnosed and fixed with an explicit
anti-copying instruction. We report this evolution because each abandoned
design is itself evidence about what does not work in this problem class,
not only as a lead-in to the final mechanism.

\subsection{The Reflective FDPO mechanism}

Reflective FDPO operationalizes the analogy from \S\ref{sec:intro}: the
optimizer is shown the measured outcome of its own previous rewrite before
it is asked to produce the next one. Algorithm~\ref{alg:reflect} gives the
precise procedure.

\begin{algorithm}[t]
\caption{Reflective FDPO}
\label{alg:reflect}
\begin{algorithmic}[1]
\Require seed prompt $p_0$, mining set $M$, validation set $V$ (disjoint
from $M$ and from the sealed test set $T$), optimizer $\textsc{Opt}$, solver
$\textsc{Solve}$, max rounds $N$
\Ensure shipped prompt $p^*$
\State $s_0 \gets \textsc{Solve}(p_0, M)$; $v_0 \gets \textsc{Solve}(p_0, V)$
\State $\text{registry} \gets \{0: (p_0, s_0, v_0)\}$
\For{$t = 1$ to $N$}
  \If{$t = 1$}
    \State $\text{reflection} \gets \emptyset$
  \Else
    \State $R_M, G_M \gets \textsc{RecoveredRegressed}(s_{t-2}, s_{t-1})$
    \State $R_V, G_V \gets \textsc{RecoveredRegressed}(v_{t-2}, v_{t-1})$
    \State $\text{reflection} \gets (R_M, G_M, |R_V|, |G_V|)$
  \EndIf
  \State $\text{failures} \gets \{i \in M : s_{t-1,i} = 0\}$
  \State $p_t \gets \textsc{Opt}(p_{t-1}, \text{failures}, \text{reflection})$
  \State $s_t \gets \textsc{Solve}(p_t, M)$; $v_t \gets \textsc{Solve}(p_t, V)$
  \State $\text{registry}[t] \gets (p_t, s_t, v_t)$ \Comment{commits every round}
\EndFor
\State $t^* \gets \arg\max_{t} \; \text{accuracy}(v_t)$
\State $p^* \gets \text{registry}[t^*].\text{prompt}$
\State \Return $p^*$ \Comment{caller evaluates $\textsc{Solve}(p^*, T)$ once}
\end{algorithmic}
\end{algorithm}

From round 2 onward, the optimizer is shown, in full and without sampling:
every mining-set item its own previous rewrite recovered or regressed, the
previous text of every section it changed, and validation-set recovery/
regression counts on a set whose item content it never otherwise sees. Every
round commits unconditionally, so a bad round never blocks a later good one
from being reachable; the final choice of which round to ship compares all
committed rounds by validation accuracy and reconstructs that round's exact
prompt from full version history, regardless of what committed afterward.
The run never reverts to the untouched seed prompt unless literally no round
ever committed.

\subsection{Evaluation metric: recovery and regression}
\label{sec:metric}

Aggregate accuracy alone cannot distinguish ``this edit changed nothing,''
``this edit fixed several items and broke none,'' and ``this edit fixed
several items while quietly breaking others'': three qualitatively
different outcomes that can share the same net accuracy delta. Let
$s_{t,i} \in \{0,1\}$ denote whether item $i$ is solved correctly under the
prompt produced at round $t$. We define, between consecutive rounds,
\begin{align}
R_t &= \{i : s_{t-1,i}=0,\, s_{t,i}=1\} \quad \text{(recovered)}\\
G_t &= \{i : s_{t-1,i}=1,\, s_{t,i}=0\} \quad \text{(regressed)}\\
N_t &= |R_t| - |G_t| \quad \text{(net transition)}
\end{align}
computed independently on the mining and validation sets every round, and
once, at the end, between the seed baseline and the shipped result on the
sealed test set. We treat this decomposition, not aggregate accuracy, as the
primary evidence behind every mechanism decision reported in this paper.

\section{Experimental Setup}
\label{sec:setup}

\begin{table*}[t]
\centering
\small
\begin{tabular}{lllccl}
\toprule
Dataset & Task & Train $\to$ mining/val & Test & Solver(s) \\
\midrule
MMLU (6 subjects) & 4-way MCQA, task-type-diverse & $50 \to 25/25$ & 66 & GPT-4o-mini, GPT-4.1 Mini \\
LegalBench-Hearsay & binary hearsay classification & $50 \to 25/25$ & 49 & GPT-4o-mini, GPT-4.1 Mini \\
IFEval/IFBench & verifiable instruction-following & $40 \to 20/20$ & 42 & GPT-4o-mini, GPT-4.1 Mini \\
AIME (2022--24 $\to$ 2025) & competition math, integer answer & $90 \to 58/32$ & 30 & GPT-4o-mini, GPT-4.1 Mini \\
PUPA & privacy-conscious delegation & $60 \to 30/30$ & 40 & GPT-4o-mini, GPT-4.1 Mini \\
\bottomrule
\end{tabular}
\caption{Experimental setup. Optimizer and judge are GPT-5 throughout;
solver, optimizer, and judge are always distinct models. Open-weight solver
runs (Qwen, Llama) are in progress; see \S\ref{sec:limitations}.}
\label{tab:setup}
\end{table*}

Table~\ref{tab:setup} summarizes the five benchmarks. All results use
Reflective FDPO with a single seed per (dataset, model) pair; we treat
single-seed results as informative but not statistically definitive (see
\S\ref{sec:limitations}).

\section{Results}
\label{sec:results}

\begin{table}[t]
\centering
\small
\begin{tabular}{llccc}
\toprule
Benchmark & Model & Base. & Final & $\Delta$ \\
\midrule
AIME & GPT-4o-mini & 13.3 & 10.0 & $-3.3$ \\
AIME & GPT-4.1 Mini & 46.7 & 53.3 & $+6.7$ \\
PUPA$^\dagger$ & GPT-4o-mini & 68.5 & 79.9 & $+11.4$ \\
PUPA$^\dagger$ & GPT-4.1 Mini & 68.2 & 80.7 & $+12.5$ \\
IFBench & GPT-4o-mini & 47.6 & 45.2 & $-2.4$ \\
IFBench & GPT-4.1 Mini & 42.9 & 66.7 & $+23.8$ \\
Hearsay & GPT-4o-mini & 71.4 & 69.4 & $-2.0$ \\
Hearsay & GPT-4.1 Mini & 71.4 & 73.5 & $+2.0$ \\
MMLU (macro) & GPT-4o-mini & 75.1 & 77.2 & $+2.0$ \\
MMLU (macro) & GPT-4.1 Mini & 81.0 & 83.3 & $+2.3$ \\
\bottomrule
\end{tabular}
\caption{Baseline vs.\ final test accuracy (\%), both solver models, all
five benchmarks. $^\dagger$PUPA reports the composite score
$(\text{quality} + (1-\text{leakage}))/2$, its paper-faithful metric. The
AIME/GPT-4.1 Mini baseline reflects a corrected, repeated re-measurement of
the seed prompt (see \S\ref{sec:limitations}); it differs from the single
originally logged run, which we no longer trust as representative.}
\label{tab:cross-benchmark}
\end{table}

Table~\ref{tab:cross-benchmark} reports baseline-vs-final test accuracy for
all five benchmarks under both solver models. Four of five benchmarks
(PUPA, IFBench, MMLU, AIME) improve under GPT-4.1 Mini; AIME regresses
slightly for GPT-4o-mini on a 30-item test set where each item is worth 3.3
percentage points; IFBench improves substantially for GPT-4.1 Mini
(+23.8pp) but regresses for GPT-4o-mini (-2.4pp) under an identical
mechanism and configuration, a genuine, disclosed model-dependent result,
not an artifact of different setups.

\subsection{Case study: PUPA and the validation/test generalization gap}

PUPA's judge role (which scores response quality) is itself a reasoning
model whose completions were, in one run, silently truncated by a hardcoded
token cap unrelated to the experiment's configured solver budget, causing
truncated judge calls to default to a quality score of zero. This was
caught only because per-item recovery/regression tracking flagged an
implausible pattern of validation-set regressions that did not correspond to
any prompt-content difference; tracing the affected items back through the
judge's raw token usage located the cap. After the fix, GPT-4.1 Mini's PUPA
composite score improved from 0.682 to 0.807 (+12.5pp), with 10 test items
recovered and zero regressed: the cleanest result in this paper, and one
that would not have been caught by aggregate accuracy alone.

\subsection{Case study: IFBench and a model-dependent regression}

IFBench's constraint checker included a specific rule
(\texttt{words:no\_consecutive}) that, on inspection of a regression the
recovery/regression trace flagged as suspicious, turned out to check a
different condition than the one stated in the actual dataset instruction
text. A second, independent issue, an optimizer-prompt instruction that had
become stale relative to the mechanism's actual round-selection logic, was
also identified through the same per-round audit trail. After fixing
both, GPT-4.1 Mini's IFBench score improved from 0.429 to 0.667 (+23.8pp,
12 recovered/2 regressed), while GPT-4o-mini's score, measured before these
fixes, remains a net regression (0.476$\to$0.452) across two identical-config
reruns. We report both numbers rather than only the favorable one.

\section{Limitations}
\label{sec:limitations}

All results in this paper use a single seed per (dataset, model)
configuration; we do not yet report variance across seeds, and at test sizes
of 30--66 items, single-item noise is 1.5--3.3 percentage points. Several
GPT-4.1 Mini results in \S\ref{sec:results} were measured after fixing bugs
discovered during this project (a judge-token-truncation bug, an IFBench
checker bug, a stale optimizer instruction); the corresponding GPT-4o-mini
results predate these fixes, so cross-model absolute-score comparisons in
Table~\ref{tab:cross-benchmark} should be read with this asymmetry in mind,
even though within-model baseline-to-final deltas are unaffected by it. We
have not yet evaluated open-weight solver models (Qwen, Llama) at a scale
directly comparable to prior work, which is necessary to answer RQ4. Our
optimizer and judge are a single frontier model (GPT-5) throughout; we have
not measured sensitivity to a weaker or different optimizer model.

% TODO(authors): expand with the Responsible NLP Research Checklist items
% required by ARR before submission.

\bibliography{custom}

\end{document}
